\documentclass[12pt]{article}
\usepackage{float}
\usepackage{graphicx}
\usepackage{fontspec}
\usepackage{url}
\usepackage{apacite}
\setmainfont{Times New Roman}
\usepackage[letterpaper,top=2cm,bottom=2cm,left=2cm,right=2cm,marginparwidth=1.5cm]{geometry}

\title{Weekly assignment 4: Reflection on social structures and their impact on organizational cybersecurity}
\author{Truong Khang Tobias Nguyen}
\date{}

\begin{document}

\maketitle

\section{Theory 1 - Critical Theory}
The first concept I chose is \textbf{critical theory}. I find the idea that knowledge, truth, and thus by extension, science is not objective to be extremely intriguing. 
\newline
The parallel that I can see between critical theory and organizational cybersecurity is that we take given security measures as objective truths to security risks. The same things goes for cybersecurity standards such as NIST CSF, ISO 27001, etc. 
\newline
We implement such standards as we believe that their measures must be correct or relevant as they are usually formulated by experts, organizations and industry. In that sense, we take their knowledge and experience for granted as we consider them the dominant group that shape our knowledge and truth as we are the oppressed group seeking guidance from them.
\newline
This is especially interesting when we consider an organization's security strategy. Typically, strategy is defined by management while the lower levels of an organization is responsible for implementing said strategy. As management develops a strategy based on past experience, knowledge and often times, advisors, one may argue that critical theory can be found in cybersecurity in a layered form where each instance bases their actions on the actions of a higher instance.
\newline
That being said, I believe that critical theory affects organizational cybersecurity in a more abstract way. The security measures and approaches we develop or follow stem either from our own experiences and knowledge or from the experts and subsequently, their opinion on security measures. This maps well to the power dynamic defined in critical theory, were objectivity is dictated by the dominant group and thus, accepted by the oppressed group.

\section{Theory 2 - Acceleration}
The other concept I chose is \textbf{acceleration}. To map this concept to organizational cybersecurity is more straightforward than critical theory which affects organizational cybersecurity in a more abstract way. 
\newline
As technology and thus, cyber threats advance, organizations also need to accelerate their efforts to address them. This means, that the rate of which security measures, approaches, strategies must continuously adapt to the advancement of technology, ideally in a linear way.
\newline
Acceleration suggests that this continuous loop ultimately leads to burnout and alienation. On an organizational level, we can describe this alienation as organizations being unable to form or adapt policies as changes happen, or to say it bluntly, can't keep up with them. As a result, organizational cybersecurity will always be behind the rate of which threats evolve. 
\newline
That being said, I believe that acceleration affects organizational cybersecurity in a way that we can understand why policies and standards become outdated in relation to technological advancements and subsequently forces organizations to become more resilient to threats.

\end{document}